\documentclass[11pt]{article}
\newcommand{\LabNumber}{1}
\newcommand{\LabTitle}{Java Fundamentals}
\newcommand{\LabTopic}{Classes, objects, OOP basics, \texttt{Scanner} I/O}
\input{common.tex}

\begin{document}
\labheader

\section*{Objectives}
\begin{itemize}[leftmargin=*,nosep]
  \item Understand the structure of a Java program: class, \texttt{main} method, packages.
  \item Declare variables of primitive and reference types; use \texttt{System.out.println}.
  \item Read user input using the \texttt{Scanner} class.
  \item Define a class with fields, a constructor, and instance methods (encapsulation).
  \item Apply the four OOP pillars: encapsulation, inheritance, polymorphism, abstraction.
\end{itemize}

\begin{summary}[Quick Reference]
\textbf{Minimal Java program.}
\begin{lstlisting}
public class Hello {
    public static void main(String[] args) {
        System.out.println("Hello, world!");
    }
}
\end{lstlisting}

\textbf{Reading input.}
\begin{lstlisting}
import java.util.Scanner;
Scanner sc = new Scanner(System.in);
int n    = sc.nextInt();
double d = sc.nextDouble();
String s = sc.next();       // one token (no spaces)
sc.close();
\end{lstlisting}

\textbf{Primitive types.}\quad \texttt{int}, \texttt{long}, \texttt{double}, \texttt{boolean}, \texttt{char}.\\
\textbf{String.}\quad Immutable reference type; use \texttt{+} for concatenation, \texttt{.length()}, \texttt{.charAt()}.\\
\textbf{Access modifiers.}\quad \texttt{public} (anywhere), \texttt{private} (class only), \texttt{protected} (class + subclasses).\\
\textbf{Workflow.}\quad Write \texttt{.java} $\to$ compile with \texttt{javac} $\to$ run with \texttt{java ClassName}.
\end{summary}

\subsection*{Quick Experiments (Try It)}

\begin{tryit}{Integer vs.\ floating-point division}
Run this program and record the output:
\begin{lstlisting}
public class DivTest {
    public static void main(String[] args) {
        System.out.println(7 / 2);
        System.out.println(7 / 2.0);
        System.out.println((double) 7 / 2);
    }
}
\end{lstlisting}
\begin{enumerate}[nosep,leftmargin=*]
  \item What is the difference between the three results?
  \item What does the cast \texttt{(double)} do?
  \item What would happen if you used \texttt{int result = 7 / 2.0;}? Try it.
\end{enumerate}
\end{tryit}

\begin{tryit}{String reference vs.\ value equality}
\begin{lstlisting}
public class StrEq {
    public static void main(String[] args) {
        String a = new String("hello");
        String b = new String("hello");
        System.out.println(a == b);
        System.out.println(a.equals(b));
    }
}
\end{lstlisting}
\begin{enumerate}[nosep,leftmargin=*]
  \item Why does \texttt{==} print \texttt{false} even though the contents are identical?
  \item Which method should always be used to compare \texttt{String} contents?
\end{enumerate}
\end{tryit}

\section*{In-Lab Exercises (target: 2 hours)}

\begin{labexercise}{Hello + Your Info \hfill {\small \itshape (15 min)}}
Write a program that prints three lines: (1) \texttt{Welcome to ADS!},
(2) your full name, (3) your student ID.
\textbf{Goal:} verify the IDE/JDK toolchain works end-to-end.
\end{labexercise}

\begin{labexercise}{GPA Calculator \hfill {\small \itshape (30 min)}}
Read four course grades (as \texttt{double}) from the user and compute and print the
average GPA formatted to two decimal places using \texttt{String.format("\%.2f", gpa)}.
Label the output clearly, e.g.\ \texttt{Your GPA: 3.56}.
\end{labexercise}

\begin{labexercise}{Defining a \texttt{Student} Class \hfill {\small \itshape (40 min)}}
Create a class \texttt{Student} with:
\begin{itemize}[leftmargin=*,nosep]
  \item \texttt{private} fields: \texttt{String name}, \texttt{int id}, \texttt{double gpa}.
  \item A constructor that initialises all three fields via \texttt{this}.
  \item Public getters for each field.
  \item An \texttt{@Override} of \texttt{toString()} that returns\\
        \texttt{Student[id=1001, name=Alice, gpa=3.80]}.
  \item A method \texttt{getLetterGrade()} returning \texttt{"A"} ($\geq 3.5$), \texttt{"B"} ($\geq 3.0$), \texttt{"C"} ($\geq 2.5$), or \texttt{"F"}.
\end{itemize}
In \texttt{main}, create two \texttt{Student} objects, print them, and print each student's letter grade.
\end{labexercise}

\begin{labexercise}{Inheritance --- \texttt{GradStudent} \hfill {\small \itshape (35 min)}}
Extend \texttt{Student} with a subclass \texttt{GradStudent} that adds:
\begin{itemize}[leftmargin=*,nosep]
  \item A field \texttt{String researchTopic}.
  \item A constructor calling \texttt{super(...)} for the inherited fields.
  \item An \texttt{@Override} of \texttt{toString()} that appends \texttt{, topic=<value>} to the parent's result.
\end{itemize}
Demonstrate polymorphism: store both a \texttt{Student} and a \texttt{GradStudent} in a
\texttt{Student} reference and call \texttt{toString()} on each.
\end{labexercise}

\begin{aicollab}
\textbf{Suggested prompts:}
\begin{itemize}[leftmargin=*,nosep]
  \item ``Explain the difference between \texttt{==} and \texttt{.equals()} for Java Strings.''
  \item ``What does \texttt{this} refer to inside a Java constructor?''
  \item ``Why does Java require \texttt{super(...)} as the first statement in a subclass constructor?''
\end{itemize}
\medskip\textbf{Verify:} Ask the AI to write a class with a private field accessed directly from
outside the class. Confirm the compiler rejects it. Note what error message appears.
\end{aicollab}

\takehomebanner

\begin{trace}[Trace \texttt{toString} output]
What does the following \texttt{main} print?
\begin{lstlisting}
public class Animal {
    String name;
    Animal(String n) { name = n; }
    public String toString() { return "Animal(" + name + ")"; }
}
public class Dog extends Animal {
    Dog(String n) { super(n); }
    public String toString() { return "Dog(" + name + ")"; }
}
// In main:
Animal a = new Dog("Rex");
System.out.println(a);
\end{lstlisting}
\end{trace}

\begin{bug}[Fix the constructor]
Find and fix all errors:
\begin{lstlisting}
public class Point {
    private int x, y;
    public Point(int x, int y) {
        x = x;   // bug
        y = y;   // bug
    }
    public String toString() {
        return "(" + x + ", " + y + ")";
    }
}
\end{lstlisting}
\end{bug}

\begin{writecode}[Rectangle class]
Write a class \texttt{Rectangle} with private fields \texttt{width} and \texttt{height} (\texttt{double}),
a constructor, an \texttt{area()} method, and a \texttt{perimeter()} method.
Write a short \texttt{main} that creates one \texttt{Rectangle} and prints both values.
\end{writecode}

\vspace{4pt}
\noindent\textbf{Submission.} Save in-lab source files as \texttt{lab1\_ex1.java}, \ldots\ and submit on the course site.
Hand in paper work at the start of the next lab.

\end{document}
